%% CS505 Final Project Proposal
%% Single-column proposal format (ACL format reserved for midway/final reports)
%% Compile: pdflatex → bibtex → pdflatex → pdflatex

\documentclass[10pt,letterpaper]{article}

%% ── Packages ──────────────────────────────────────────────────────────────────
\usepackage[T1]{fontenc}
\usepackage[utf8]{inputenc}
\usepackage{times}
\usepackage[top=0.9in, bottom=0.9in, left=1in, right=1in]{geometry}
\usepackage{natbib}
\usepackage{microtype}
\usepackage{booktabs}
\usepackage{amsmath}
\usepackage{amssymb}
\usepackage{url}
\usepackage[hidelinks]{hyperref}
\usepackage{parskip}

\makeatletter
\renewcommand{\section}{\@startsection{section}{1}{\z@}%
  {7pt \@plus 1pt \@minus 1pt}{2pt}%
  {\normalfont\normalsize\bfseries}}
\renewcommand{\subsection}{\@startsection{subsection}{2}{\z@}%
  {4pt \@plus 1pt}{1.5pt}%
  {\normalfont\normalsize\itshape}}
\makeatother
\setlength{\parskip}{3.5pt}
\setlength{\parindent}{0pt}
\bibsep=2.5pt

%% ── Title ─────────────────────────────────────────────────────────────────────
\title{\vspace{-1.5em}%
  \textbf{Reducing Hallucinations via Citation-Grounded QA
  with Locally-Hosted Small Language Models}\\[3pt]
  {\normalsize\normalfont
   CS505 Natural Language Processing \enspace|\enspace
   Final Project Proposal \enspace|\enspace Boston University}\\[1pt]
  {\normalsize\normalfont [We will add names here :)]}%
}
\date{\vspace{-1.5em}}

\begin{document}
\maketitle
\thispagestyle{empty}

%% ── Preamble ─────────────────────────────────────────────────────────────────
\noindent\textbf{Project type:}
Empirical NLP project focused on retrieval-augmented generation (RAG) for
hallucination reduction.

\noindent\textbf{Research question:}
\emph{How much of the hallucination gap between a locally-hosted small language
model (SLM) and GPT-4o is attributable to retrieval quality versus generation
capacity—and can systematic retrieval improvements make fully local,
privacy-preserving deployment practically viable?}

%% ── 1. Problem Statement ─────────────────────────────────────────────────────
\section{Problem Statement \& Motivation}

Hallucination—generating claims unsupported by source documents—is the
primary failure mode of LLMs in high-stakes settings.
Cloud tools like Google's NotebookLM answer questions over private documents
with inline citations, but require transmitting the entire corpus to an
external server, violating HIPAA, GDPR, and corporate data-governance
policies.
\textbf{Cloud upload equals data-exfiltration risk.}
Sub-7B SLMs (e.g., Phi-3-mini-3.8B \citep{microsoft2024phi3}) can run on a
single consumer GPU, but it is unclear whether their higher hallucination rate
stems from \emph{retrieval failures} or \emph{generation failures}.
If retrieval improvements close most of the gap, local deployment becomes a
practical alternative.

\textbf{Task definition.}
Given a corpus $\mathcal{C}=\{d_1,\ldots,d_n\}$ and query $q$, the system
produces answer $a$ plus, for every sentence $s_i\in a$, a citation
$(p_i, s_i)$ where passage $p_i\in\mathcal{C}$ textually entails $s_i$.
Any sentence without a valid citation is counted as a hallucination.

%% ── 2. Proposed System ───────────────────────────────────────────────────────
\section{Proposed System}

All components run locally (no API calls except the GPT-4o oracle).

\textbf{Parsing.}  PyMuPDF; preserves section headings for structure-aware
chunking.

\textbf{Chunking (ablated).}  Three strategies: fixed-size (256 tok, 32
overlap), sliding-window (512 tok, 128 overlap), section-aware
(boundary-aligned).

\textbf{Retrieval (ablated).}  Dense-only (FAISS + BGE-small-en-v1.5
\citep{xiao2023bge}); BM25 (rank-bm25); hybrid reciprocal-rank fusion (RRF);
hybrid + cross-encoder reranking (ms-marco-MiniLM-L-6-v2), top-5 passages.

\textbf{Generation.}  Two prompt variants: standard generation (no citation
instruction) vs.\ \emph{citation-forcing} (the model must bracket every claim
with an inline citation and assert only what the retrieved passages support).
Served via Ollama; $\leq$16 GB VRAM.

%% ── 3. Experimental Design ───────────────────────────────────────────────────
\section{Experimental Design}

\subsection{Conditions}

Each condition changes exactly one component (Table~\ref{tab:ablation}).
Conditions A--E are the committed core; F and G are stretch goals.

\begin{table}[h]
\centering
\small
\setlength{\tabcolsep}{5pt}
\renewcommand{\arraystretch}{1.1}
\begin{tabular}{@{}clllcc@{}}
\toprule
\textbf{Cond.} & \textbf{Chunk} & \textbf{Retrieval} &
\textbf{Rerank} & \textbf{Cite-force} & \textbf{Generator} \\
\midrule
A (base) & Fixed   & Dense  & $\times$   & $\times$   & Phi-3-mini \\
B        & Fixed   & Hybrid & $\times$   & $\times$   & Phi-3-mini \\
C        & Section & Hybrid & $\times$   & $\times$   & Phi-3-mini \\
D        & Section & Hybrid & \checkmark & $\times$   & Phi-3-mini \\
E (best) & Section & Hybrid & \checkmark & \checkmark & Phi-3-mini \\
\midrule
F$^*$    & Section & Hybrid & \checkmark & \checkmark & Llama-3-8B \\
G$^*$    & Section & Hybrid & \checkmark & \checkmark & GPT-4o (oracle) \\
\bottomrule
\end{tabular}
\caption{Ablation conditions. $^*$Stretch goals (time permitting).}
\label{tab:ablation}
\end{table}

\subsection{Dataset \& Baselines}

\textbf{Dataset:} QASPER \citep{dasigi2021qasper}—1,585 full-text NLP papers,
5,049 Q\&A pairs with gold evidence annotations.
Supports both retrieval evaluation and answer evaluation (official token F1).

\textbf{Baselines:}
(1)~\textbf{Parametric-only LM}—Phi-3-mini with no retrieval, establishing
the hallucination floor.
(2)~\textbf{BM25-only}—sparse keyword retrieval, the simplest retrieval
baseline.

\subsection{Metrics \& Hypothesis}

Retrieval: Recall@5, MRR@10, evidence hit rate.
Answer: Token F1 (QASPER official).
Faithfulness: RAGAS \citep{es2023ragas}.
Citation precision: NLI entailment (TRUE \citep{honovich2022true}).

\textbf{Hypothesis:} Retrieval improvements (A$\to$D) will reduce hallucination
more than citation-forcing alone (D$\to$E), identifying retrieval as the
primary bottleneck; section-aware chunking will boost citation precision.

%% ── 4. Related Work ──────────────────────────────────────────────────────────
\section{Related Work}

\citet{lewis2020rag} introduced RAG; \citet{rashkin2023ais} and
\citet{gao2023alce} defined attribution evaluation.
RAGAS \citep{es2023ragas} automates faithfulness scoring.
\citet{ajith2024litsearch} motivates hybrid retrieval for scientific text;
\citet{asai2024openscholar} is the closest prior system.
Our contribution—a local-SLM ablation disentangling retrieval from generation
failures—is absent from all prior work.

%% ── 5. Feasibility & Timeline ────────────────────────────────────────────────
\section{Feasibility \& Timeline}

\textbf{Scope.}
The committed core (A--E on QASPER with Phi-3-mini) builds the retrieval
pipeline incrementally—one component per week—and is realistic on a single
consumer GPU.
Stretch goals (F, G, embedding ablation, human eval) will be attempted in
week 5 and deferred to the final report if time is short.

\begin{table}[h]
\centering
\small
\setlength{\tabcolsep}{5pt}
\renewcommand{\arraystretch}{1.05}
\begin{tabular}{@{}cp{9cm}@{}}
\toprule
\textbf{Wk} & \textbf{Milestone} \\
\midrule
1--2 & Parsing, chunking strategies, BM25 + dense baseline (cond.\ A) \\
3    & Hybrid RRF + reranking; conditions B--D; QASPER eval pipeline \\
4    & Citation-forcing (E); RAGAS + NLI evaluation; error analysis \\
5    & Midway report; begin stretch goals (F, G) if on schedule \\
6    & Final ablation, writeup (ACL format), GitHub release (MIT) \\
\bottomrule
\end{tabular}
\caption{Six-week schedule. Core scope: weeks 1--4 (conditions A--E).}
\end{table}

\noindent\textbf{Stack:} PyMuPDF, sentence-transformers, FAISS, rank-bm25,
Ollama, RAGAS, HuggingFace Transformers; MIT license; no paid API for core scope.

%% ── References ───────────────────────────────────────────────────────────────
{\footnotesize
\bibliographystyle{abbrvnat}
\bibliography{references}
}

\end{document}
